\reading{IM-2}{Factors}{STRIKE FIRST}{5}

\srcnote{Andrew Ang, \textit{Asset Management: A Systematic Approach to Factor
Investing} (New York, NY: Oxford University Press, 2014), Chapter 7.
Printed as Chapter 2 in the GARP source volume.}

\begin{imkeybox}[title={The five objectives in this reading}]
\begin{enumerate}[label=\textbf{\alph*.}]
\item Describe the process of value investing and explain why a value premium may exist.
\item Explain how different macroeconomic risk factors, including economic growth, inflation, and volatility, affect asset returns and risk premiums.
\item Assess methods of mitigating volatility risk in a portfolio and describe challenges that arise when managing volatility risk.
\item Explain how dynamic risk factors can be used in a multifactor model of asset returns, using the Fama-French model as an example.
\item Compare value and momentum investment strategies, including their return and risk profiles.
\end{enumerate}
\end{imkeybox}

\basics{Basics: the two types of factor}

There are two types of factors.

\begin{itemize}
\item The first type is \term{macro, fundamental-based factors}, which include
economic growth, inflation, volatility, etc.
\item The second type is \term{investment-style factors} like the market factor
of the capital asset pricing model (CAPM). To be more specific, investment-style
factors can be divided into:
  \begin{itemize}
  \item \term{Static factors}, like the market factor in CAPM, which we simply go
  long to collect a risk premium;
  \item \term{Dynamic factors}, which can only be exploited through constantly
  trading different types of securities.
  \end{itemize}
\end{itemize}

% ------------------------------------------------------------------
\lo{IM-2 a}{Describe the process of value investing and explain why a value premium may exist.}

\subsection{Value investing}

\begin{imfmlbox}
\[ \text{Book-to-market ratio} \;=\; \frac{\text{book value of equity}}{\text{market capitalisation}}
   \;=\; \frac{\text{total assets} - \text{total liabilities}}{\text{share price} \times \text{shares outstanding}} \]
\vk{book value of equity is total assets minus total liabilities; market
capitalisation is share price times shares outstanding. Equivalently, book value
per share divided by price per share.}
\end{imfmlbox}

\begin{imtrapbox}[title={A missing denominator in the source notes}]
The notes give the book-to-market ratio as ``total assets minus total
liabilities / shares outstanding.'' That expression is \term{book value per
share}, which is the numerator only, and it returns a currency amount rather than
a ratio. GARP's definition is explicit: \emph{the book-to-market ratio is book
value divided by market capitalisation}. Without the price in the denominator
there is no ratio, and no way to sort stocks into value and growth --- which is
the entire mechanism of this objective.
\end{imtrapbox}

\begin{itemize}
\item High ratio $=$ \term{value stock} (potentially undervalued).
\item Low ratio $=$ \term{growth stock} (potentially overvalued).
\end{itemize}

\term{Value investing strategy:} buying value stocks (high book-to-market) and
shorting growth stocks (low book-to-market).

\term{Growth stocks investing:} buy low book-to-market ratio stocks.

\subsection{The value premium}

Historically, value stocks outperform growth stocks (but with periods of
underperformance). This outperformance may be due to a value risk. Investors are
compensated for:

\begin{itemize}
\item \term{Rational explanation:} value stocks are inherently riskier, demanding
a higher return.
\item \term{Behavioral explanation:} investors underestimate value stocks and
overestimate growth stocks, creating buying opportunities in undervalued
companies.
\end{itemize}

% ------------------------------------------------------------------
\lo{IM-2 b}{Explain how different macroeconomic risk factors, including economic growth, inflation, and volatility, affect asset returns and risk premiums.}

\subsection{Macroeconomic factors}

\renewcommand{\arraystretch}{1.25}
\noindent\begin{tabularx}{\textwidth}{@{}>{\raggedright\arraybackslash}p{3.0cm}X@{}}
\toprule
\textbf{\sffamily\small Factor} & \textbf{\sffamily\small What the notes say happens} \\
\midrule
Economic growth & Risky assets generally perform poorly and are much more volatile during periods of \emph{low} economic growth. However, government bonds tend to do well during these times. \\
Inflation & High inflation tends to be bad for \emph{both} stocks and bonds. During periods of high inflation, all assets tend to do poorly. Part of the long-run risk premiums for both equities and bonds represents compensation for doing badly when inflation is high. \\
Productivity risk & When productivity slows, stock returns tend to be low. \\
Demographic risk & Labor income is mostly earned and saved in young and middle age and dis-saved when retired. \\
Political risk (sovereign risk) & Mainly in emerging countries. In financial crisis, developed countries have political risk as well. \\
Volatility & The negative relation between volatility and stock returns is called the \term{leverage effect}. \\
\bottomrule
\end{tabularx}
\renewcommand{\arraystretch}{1}
\figcap{Six macro factors and their direction of effect. Growth is the only row
where stocks and bonds move in opposite directions; inflation is the row where
they move together, and together downwards. That contrast is the examinable
distinction.}

\subsection{Volatility}

There are two channels running from volatility to stock returns:

\begin{enumerate}
\item \term{The leverage effect.} When stock returns drop, the financial leverage
of firms increases since debt is approximately constant while the market value of
equity has fallen. This makes equities riskier and increases their volatilities.
\item \term{The required-return channel.} An increase in volatility raises the
required return on equity demanded by investors, also leading to a decline in
stock prices.
\end{enumerate}

Bonds offer some but not much respite during periods of high volatility, as the
correlation between bond returns and VIX changes is only $0.12$. Thus, bonds are
not always a safe haven when volatility shocks hit.

\figwrap{%
\begin{tikzpicture}
\begin{axis}[frmaxis,
  width=11.4cm, height=5.2cm,
  ybar, bar width=26pt,
  symbolic x coords={VIX vs stocks, VIX vs bonds, stocks vs bonds},
  xtick=data, x tick label style={font=\scriptsize\sffamily},
  ymin=-0.58, ymax=0.30,
  ytick={-0.4,-0.2,0,0.2},
  ylabel={correlation of monthly changes},
]
\addplot[draw=FRMRed,fill=PaleRed,area legend] coordinates
  {(VIX vs stocks,-0.39) (VIX vs bonds,0.12) (stocks vs bonds,-0.01)};
\draw[black!45, thin] (axis cs:VIX vs stocks,0) -- (axis cs:stocks vs bonds,0);
\node[font=\scriptsize\sffamily,color=black!70,anchor=north,fill=white,inner sep=1.5pt]
  at (axis cs:VIX vs stocks,-0.42) {$-0.39$};
\node[font=\scriptsize\sffamily,color=black!70,anchor=south,fill=white,inner sep=1.5pt]
  at (axis cs:VIX vs bonds,0.14) {$0.12$};
\node[font=\scriptsize\sffamily,color=black!70,anchor=north,fill=white,inner sep=1.5pt]
  at (axis cs:stocks vs bonds,-0.05) {$-0.01$};
\end{axis}
\end{tikzpicture}}%
{The three correlations behind the safe-haven claim, monthly, March 1986 to
December 2011 (GARP source, Chapter 2). Stocks fall hard when volatility rises,
at $-0.39$. Bonds rise, but only at $0.12$ --- a fifth of the magnitude, and the
point of the sentence above. The near-zero stock--bond figure of $-0.01$ is the
one most often misremembered as strongly negative.}

% ------------------------------------------------------------------
\lo{IM-2 c}{Assess methods of mitigating volatility risk in a portfolio and describe challenges that arise when managing volatility risk.}

\subsection{Mitigating volatility risk --- buying protection}

Investors who dislike volatility risk can \term{buy volatility protection}.
Buying volatility protection can benefit from increasing market volatility. There
are many ways to buy volatility protection:

\begin{itemize}
\item In option markets, but traders can also use other derivatives contracts,
such as volatility swaps.
\item Assets that pay off during high volatility periods, like out-of-the-money
puts, provide hedges against volatility risk.
\item Investors allergic to volatility could increase their holdings of bonds,
but bonds do not always pay off during highly volatile periods.
\end{itemize}

\subsection{Investing in volatility risk --- selling protection}

Investors who love volatility risk can \term{sell volatility protection}.
Although selling volatility produces high and steady payoffs during stable times,
once every decade or so there is a huge crash where sellers of volatility
experience large, negative payoffs. Unfortunately, some investors who sold
volatility prior to the financial crisis failed to anticipate that a crash like
the one of 2008 would materialize.

\figwrap{%
\begin{tikzpicture}
\begin{axis}[frmaxis,
  width=11.6cm, height=5.4cm,
  xlabel={time}, ylabel={cumulative payoff},
  xmin=0, xmax=10.2, ymin=-6.6, ymax=3.6,
  xtick=\empty,
  legend style={font=\scriptsize\sffamily, draw=FRMRule, fill=white,
    at={(0.5,-0.24)}, anchor=north, legend columns=-1, column sep=10pt},
]
\addplot[FRMGreen, very thick, sharp plot] coordinates
  {(0,0) (1,0.35) (2,0.70) (3,1.05) (4,1.40) (5,1.75) (6,2.10) (7,2.45)
   (7.5,2.60) (7.8,-5.4) (8.2,-5.0) (9,-4.6) (10,-4.2)};
\addlegendentry{selling volatility protection}
\addplot[FRMNavy, very thick, dashed, sharp plot] coordinates
  {(0,0) (1,-0.30) (2,-0.60) (3,-0.90) (4,-1.20) (5,-1.50) (6,-1.80) (7,-2.10)
   (7.5,-2.25) (7.8,2.6) (8.2,2.4) (9,2.1) (10,1.8)};
\addlegendentry{buying volatility protection}
\draw[black!45, thin] (axis cs:0,0) -- (axis cs:10.2,0);
\node[font=\scriptsize\sffamily, fill=white, inner sep=1.5pt, anchor=west,
      text=FRMRed] (cr) at (axis cs:3.2,-5.2) {the once-a-decade crash};
\draw[-{Stealth[length=4pt]}, FRMRed, thin] (cr.east) -- (axis cs:7.55,-4.6);
\end{axis}
\end{tikzpicture}}%
{The shape of the two positions, and why the choice is not symmetric in
experience even though it is symmetric in construction. Selling protection is a
long run of small steady gains ending in one very large loss; buying protection
is a long run of small steady losses ending in one very large gain. Both lines
are the same trade seen from opposite sides. The seller's record looks better for
years, which is exactly the difficulty.}

\subsection{The challenge in managing volatility risk}

Constructing valuation models with volatility risk can be tricky because the
relation between volatility and expected returns is \term{time varying and
switches signs}, and is thus very hard to pin down.

\begin{imtrapbox}[title={Direction is the whole question here}]
Buy protection if you \emph{dislike} volatility risk; sell protection if you
\emph{love} it. The payoff asymmetry runs the other way from the intuition: the
seller earns steadily and loses catastrophically, not the reverse. And the
modelling problem is not that the sign is unknown, but that it \emph{changes}
over time.
\end{imtrapbox}

% ------------------------------------------------------------------
\lo{IM-2 d}{Explain how dynamic risk factors can be used in a multifactor model of asset returns, using the Fama-French model as an example.}

The best-known example of a tradeable multifactor model applied to present
dynamic factors is introduced by Fama and French.

The Fama--French (1993) model explains asset returns with three factors. There is
the traditional CAPM market factor and there are two additional factors to
capture a size effect and a value/growth effect:

\begin{imfmlbox}
\[ E(r_i) = r_f + \beta_{i,\text{MKT}}E(r_m - r_f) + \beta_{i,\text{SMB}}E(\text{SMB})
   + \beta_{i,\text{HML}}E(\text{HML}) \]
\vk{$\beta_{i,\text{MKT}}$, $\beta_{i,\text{SMB}}$ and $\beta_{i,\text{HML}}$ are
the asset's loadings on the three factors, and $E(\cdot)$ denotes each factor's
risk premium.}
\end{imfmlbox}

\begin{itemize}
\item \term{SMB factor}, which refers to (buy) small stocks minus (sell) big
stocks. Small and big refer to market capitalisation.
\item \term{HML factor}, which stands for (buy) high book-to-market stocks minus
a portfolio of (sell) low book-to-market stocks.
\end{itemize}

\begin{imkeybox}
Both factors are \term{long-short}. That is what makes them dynamic rather than
static: the portfolio has to be re-sorted and re-traded as firms change size and
valuation, whereas the market factor can be held passively.
\end{imkeybox}

% ------------------------------------------------------------------
\lo{IM-2 e}{Compare value and momentum investment strategies, including their return and risk profiles.}

\subsection{Momentum}

Momentum is the strategy of buying stocks that have gone up over the past six (or
so) months (\term{winners}) and shorting stocks with the lowest returns over the
same period (\term{losers}).

In 1993, Jegadeesh and Titman identified a momentum effect, which is then
introduced by Carhart in his Four-Factor model. The momentum effect refers to the
phenomenon that winner stocks continue to win and losers continue to lose, just
like the ``Matthew effect.''

\begin{imtrapbox}[title={The factor has a name, and the notes omit it}]
The momentum factor is \term{WML} --- past \emph{winners minus losers} --- also
written \term{UMD}, for stocks that have gone \emph{up minus down}. It sits
alongside SMB and HML as a named factor and is examinable as one. The source
notes give the strategy without ever naming the factor. Also note the spelling:
GARP writes \term{Jegadeesh}, not ``Jagadeesh.''
\end{imtrapbox}

\subsection{Comparison of the two strategies}

\begin{imdefbox}[title={The same}]
The momentum strategy, like size and value, is a \term{cross-sectional} strategy,
meaning that it compares one group of stocks against another group of stocks in
the cross section, rather than looking at a single stock over time.
\end{imdefbox}

\figwrap{%
\begin{tikzpicture}
\begin{axis}[frmaxis, name=left, scale only axis,
  width=5.6cm, height=4.2cm,
  title={Value: negative feedback},
  xlabel={time}, ylabel={price},
  xmin=0, xmax=10, ymin=42, ymax=118, xtick=\empty,
]
\addplot[FRMNavy, very thick, smooth] coordinates
  {(0,100) (1,92) (2,84) (3,74) (4,64) (5,58) (6,66) (7,76) (8,85) (9,92) (10,97)};
\addplot[only marks, mark=*, mark size=2pt, FRMGreen] coordinates {(5,58)};
\node[font=\scriptsize\sffamily, fill=white, inner sep=1.5pt, anchor=north west,
      text=FRMGreen] (bv) at (axis cs:0.3,114) {buy \emph{after} the fall};
\draw[-{Stealth[length=4pt]}, FRMGreen, thin] (bv.south) -- (axis cs:4.7,61);
\end{axis}
\begin{axis}[frmaxis, name=right, at={(left.right of south east)}, xshift=1.9cm,
  scale only axis, width=5.6cm, height=4.2cm,
  title={Momentum: positive feedback},
  xlabel={time},
  xmin=0, xmax=10, ymin=42, ymax=118, xtick=\empty, ytick=\empty,
]
\addplot[FRMNavy, very thick, smooth] coordinates
  {(0,58) (1,64) (2,71) (3,79) (4,86) (5,92) (6,97) (7,101) (8,104) (9,106) (10,107)};
\addplot[only marks, mark=*, mark size=2pt, FRMRed] coordinates {(5,92)};
\node[font=\scriptsize\sffamily, fill=white, inner sep=1.5pt, anchor=north west,
      text=FRMRed] (bm) at (axis cs:0.3,114) {buy \emph{because} of the rise};
\draw[-{Stealth[length=4pt]}, FRMRed, thin] (bm.south east) -- (axis cs:4.8,90);
\end{axis}
\end{tikzpicture}}%
{The same buy decision, triggered by opposite signals, with opposite consequences
for the price. The value buyer arrives when the price has already fallen far
enough to make the stock cheap relative to book, and the buying pushes it back
towards fundamentals --- a \emph{stabilising} trade. The momentum buyer arrives
because the price is rising and buys more of what has already risen,
pushing it further from fundamentals --- a \emph{destabilising} trade.}

\begin{imdefbox}[title={The difference}]
\begin{itemize}
\item \term{Value is a negative feedback strategy} (stabilising): stocks with
declining prices eventually fall far enough that they become value stocks.
\item \term{Momentum is a positive feedback strategy} (destabilising): stocks
with high past returns are attractive, momentum investors continue buying them,
and they continue to go up.
\end{itemize}
\end{imdefbox}

\begin{imkeybox}[title={Remember}]
The cumulated profits of momentum have been significantly larger than those of
size or value strategies.
\end{imkeybox}

\begin{imtrapbox}[title={Consolidated trap summary --- IM-2}]
\begin{itemize}
\item \term{Formula, IM-2 a.} Book-to-market needs market capitalisation in the
denominator. Book value per share alone is not the ratio.
\item \term{Polarity, IM-2 a.} \emph{High} book-to-market is value; \emph{low} is
growth. The strategy is long high, short low.
\item \term{Polarity, IM-2 b.} Low growth hurts risky assets and helps government
bonds. High inflation hurts both. Only one of these two rows is a hedge.
\item \term{Magnitude, IM-2 b.} Bond returns and VIX changes correlate at $0.12$,
not at some large positive number. Bonds are only a partial haven.
\item \term{Sibling, IM-2 b.} The \emph{leverage effect} is the channel running
from falling equity value to rising leverage to higher volatility. The
required-return channel is a separate, second mechanism. Do not merge them.
\item \term{Direction, IM-2 c.} Dislike volatility $\Rightarrow$ buy protection.
Love volatility $\Rightarrow$ sell protection. The seller's loss is rare and
large.
\item \term{Definition, IM-2 d.} SMB is small minus big by market capitalisation;
HML is high minus low by book-to-market. Swapping the sort variable between them
is the obvious corruption.
\item \term{Sibling, IM-2 e.} Value is negative feedback and stabilising;
momentum is positive feedback and destabilising. Both are cross-sectional, and
that shared property is where a distractor will claim a difference.
\end{itemize}
\end{imtrapbox}


% ==================================================================
